% Generated draft: requires mathematical and source review.
\documentclass[11pt]{article}
\usepackage[margin=25mm]{geometry}
\usepackage{fontspec}
\setmainfont{Latin Modern Roman}
\usepackage{amsmath,amssymb,mathtools,unicode-math}
\setmathfont{Latin Modern Math}
\usepackage{enumitem,needspace,xurl,hyperref}
\setlength{\parindent}{0pt}
\setlength{\parskip}{5pt}
\setlength{\emergencystretch}{3em}
\newcommand{\N}{\mathbb{N}}
\newcommand{\Z}{\mathbb{Z}}
\newcommand{\Q}{\mathbb{Q}}
\newcommand{\R}{\mathbb{R}}
\newcommand{\C}{\mathbb{C}}
\newcommand{\F}{\mathbb{F}}
\newcommand{\A}{\mathbb{A}}
\newcommand{\PP}{\mathbb{P}}
\newcommand{\E}{\mathbb{E}}
\newcommand{\eps}{\varepsilon}
\newcommand{\dd}{\,\mathrm{d}}
\begin{document}
\section*{UnsolvedMath ID: 30005601}
\textbf{Source problem code:} catalogue entry
\subsection{Definitions and mathematical statement}
Under the standard Euclidean interpretation of the supplied statement, let \(n\geq 1\) and let \(\Omega\subset\R^n\) be a nonempty bounded open convex set. Write
\[
0=\mu_0(\Omega)<\mu_1(\Omega)\leq\mu_2(\Omega)\leq\cdots
\]
for the eigenvalues, counted with multiplicity, of the nonnegative Neumann Laplacian. The boundary condition is understood weakly: an eigenfunction \(u\in H^1(\Omega)\setminus\{0\}\) with eigenvalue \(\mu\) satisfies
\[
\int_\Omega \nabla u\cdot\nabla v\,\dd x
=\mu\int_\Omega uv\,\dd x
\qquad\text{for every }v\in H^1(\Omega).
\]
This interpretation imposes no additional smoothness assumption on the boundary.

The dimension-free adjacent-eigenvalue conjecture is the assertion
\[
\exists C\in(0,\infty)\quad
\forall n\geq1\quad
\forall\,\Omega\subset\R^n\text{ nonempty, bounded, open and convex}\quad
\forall k\geq1:\qquad
\mu_{k+1}(\Omega)\leq C\mu_k(\Omega).
\]
Here \(C\) is independent of the dimension, the domain, and the eigenvalue index. The source additionally asks for an analytic proof. The displayed indexing is an explicit convention for the positive spectrum: the source's phrase ``any \(k\)'' cannot include comparison of the first positive eigenvalue with the zero eigenvalue.

The curated statement also requests an extension to convex domains in Riemannian manifolds of nonnegative Ricci curvature. Its intended form appears to be an analogous universal bound
\[
\mu_{k+1}(\Omega,g)\leq C\mu_k(\Omega,g),\qquad k\geq1,
\]
for Neumann eigenvalues, uniformly over an intended class of ambient manifolds and convex domains. That class is not specified sufficiently to make this extension a fully quantified conjecture: the meaning of convexity, completeness or compactness hypotheses, boundary assumptions, conditions ensuring a discrete Neumann spectrum, and permitted dependence of \(C\) remain unresolved.

The verbatim extract also requests an analytic proof of an inequality labelled (4), but neither (4) nor the referenced inequality (1) is supplied. Consequently, no separate mathematical statement for those inequalities can be reconstructed from this record.
\subsection{Short English statement}
The main conjecture asks whether successive positive Neumann eigenvalues of every bounded convex Euclidean domain have a ratio bounded by one constant independent of dimension, domain, and index, and asks for an analytic proof. A requested extension to manifolds with nonnegative Ricci curvature lacks necessary formulation details. The supplied open status and verification labels are curation claims requiring human review.
\subsection{Sources}
\begin{itemize}
\item Local UnsolvedMath registry, permanent ID \(30005601\), source problem code \texttt{catalogue entry}, title: Dimension-Free Neumann Eigenvalue-Ratio Bounds. This draft uses the supplied complete registry record; repository registry locations are \texttt{data/unsolvedmath/problems.json} and \texttt{catalog/problems.json}.
\item Supplied report citation: \emph{Geometric Spectral Theory (2024)}, Oberwolfach Reports, pp.\ 2017--2106, DOI \texttt{10.4171/owr/2023/36}. The literature-source label is ``Oberwolfach Report 36'': \url{https://doi.org/10.4171/owr/2023/36}.
\item Supplied literature source labelled ``Lagace primary report'': \url{https://nms.kcl.ac.uk/jean.lagace/papers/Hom_Control.pdf}. No further bibliographic identification or attribution of this conjecture to a named person is supplied.
\item The original extract attributes the conjecture to its unnamed speaker and cites references [7, 8], whose bibliographic entries are absent. It refers to inequalities (1), (3), and (4); only the adjacent-eigenvalue comparison identified as (3) is represented in the supplied extract.
\item Literature assessment dated exactly \texttt{2026-08-22}: the supplied curation labels the problem open and states that known quantitative Neumann eigenvalue inequalities for convex domains do not establish the conjectured adjacent-eigenvalue bound, and that the requested analytic proof remains open. This is supplied dated curation, not an independently verified literature assessment.
\item The record labels its statement \texttt{corrected\_verified} and says it was reconstructed and checked against the official report. These are supplied provenance claims; no independent source verification is asserted here. The record gives proposed year 2023, whereas its report citation includes 2024; both are retained without resolving their relationship.
\end{itemize}
\end{document}
